
\section{Preliminaries}\label{sec:prelim}
In this section, we introduce some useful notations and fundamental concepts of the FJ model of opinion formation for the convenience of description of problem formulation and algorithm.

\subsection{Notations}
We use $\mathbb{R}$ to represent the set of real numbers and regular lowercase letters like $a,b,c$ to represent scalars in $\mathbb{R}$. Bold lower letters are column vectors, e.g., $\aaa,\bb,\cc$. Bold capital letters, e.g., $\AA,\BB,\CC$ are matrices. To denote specific elements in vectors and matrices, We denote $A_{ij}$ as the entry of $\AA$ at $i$-th row and $j$-th column and $a_i$ as $i$-th element of vector $\aaa$. And we use $\aaa^\top$ and $\AA^\top$ to denote transpose of vector $\aaa$ and matrix $\AA$. 
We use $\one$ to denote the vector of appropriate dimensions with all entries being ones and use $\ee_u$ to represent the $u$-th standard basis vector of appropriate dimensions. 
$\II$ denotes the identify matrix of appropriate dimensions. 

\subsection{Graph and Matrix Definitions}
Let $\calG = (V,E)$ be an unweighted simple graph with $n = |V|$ nodes and  $m = |E|$ edges, where $V$ is the set of nodes and $E \subseteq V \times V$ is the set of edges. The graph $\calG$ is either a directed graph or an undirected graph depending on the context. The adjacency relation of all nodes in $\calG$ is encoded in the adjacency matrix $\AA$,  whose entry $a_{i j}=1$ if $i$ and $j$ are adjacent, and $a_{i j}=0$ otherwise. For any given node $i$, $N_i$ denote the  set  of its neighbors, meaning $N_i =\{ j: (i,j)\in E\}$. The degree \(d_i\) of any node \(i\) is defined by \(d_i=\sum_{j=1}^n a_{ij}\). The degree diagonal matrix of $\calG$ is $\DD={\rm diag}(d_1, d_2, \cdots, d_n)$. The degree vector $\dd$ is given by $\dd = \DD\one$. The transition matrix of an unbiased random walk on graph $\calG$ is $\PP = \DD^{-1} \AA$, which is a row-stochastic matrix. The Laplacian matrix $\LL$ of $\calG$ is $\LL = \DD - \AA$, which is a symmetric semi-positive definite matrix. Matrix $\LL$ is singular, and thus it is irreversible. For a matrix $\MM$, if the magnitudes of all the eigenvalues of $M$ are less than $1$, then the series $\sum_{k=0}^\infty \MM^k$ is a Neumann series and converges to \( (\II - \MM)^{-1} \). 

\subsection{Friedkin--Johnsen (FJ) Model}\label{sec:FJ}
As one of the first opinion dynamics models, the FJ model is an extension of the DeGroot's opinion model. In the DeGroot
model, every node has only one opinion that is updated as the weighted
average of its neighbors. Different from the DeGroot model, in the FJ
model, each node $i \in V$ has two different kinds of opinions: one is
the internal (or innate) opinion $s_i$, the other is the expressed
opinion $z_i$. The internal opinion $s_i$ is assumed to remain constant,
private to node $i$, while the expressed opinion $z_i$ evolves as a
weighted average of its corresponding internal opinion $s_i$ and the
expressed opinions of $i$'s neighbors. 


Given a graph $G=(V,E)$, let $\mathbf{s}$ denote the internal opinion vector, with each entry satisfying $s_i\in[0,1]$.
The formulation of the equilibrium opinion expressed in the FJ model is an iterative form such that
\begin{equation}
z^{(t+1)}_{i}=\frac{s_{i}+\sum_{j\in N_{i}}a_{ij}z^{(t)}_{j}}{1+\sum_{j\in N{i}}a_{ij}}.
\label{fj-iteration}
\end{equation}

As in the FJ iterative formula above, the opinion of each node in the FJ model is updated as the weighted average of the neighbor's expressed opinion $z_j$ and its own internal opinion $s_i$. 
When $t$ approches infinity, the updated expressed opinion was shown to converge to a unique equilibrium opinion $\zz_i$. So the iterative update in Equation~1 can be written compactly as 
\(
{\zz}^{(t+1)} = (\II+\DD)^{-1}\bigl(\s + \AA\zz^{(t)}\bigr)
\)
in vector form,
which expresses the FJ dynamics as a linear transformation applied at each iteration.
Then the system of equilibrium
expressed opinion vector $\zz$ is equivalent to:
$(\II+\LL)\zz=\s,$
where $\II$ is the identify matrix , and $\LL=\DD-\AA$ is the Laplacian matrix. When $\mathcal{G}$ is connected and $a_{ij}\ge 0$, the matrix $\II+\LL$ is symmetric positive definite, ensuring existence and uniqueness of $\zz$. 
This characterization motivates estimators and algorithms that exploit the spectral and structural properties of the Laplacian. 

\subsection{Game-theoretic Interpretation}
From a game-theoretic perspective, the classical FJ model admits an equivalent interpretation as the equilibrium of a quadratic utility game.
Recall that the FJ equilibrium satisfies 
$(\II+\DD-\AA)\zz=\s$.
For undirected graphs, the matrix $\II+\DD-\AA$ corresponds to a symmetric, positive definite Laplacian-based operator, which allows the equilibrium to be characterized as the unique minimizer of a strongly convex quadratic objective.

This formulation further admits a decentralized interpretation.
Updating $z_i$ according to the FJ rule is equivalent to each agent $i$
minimizing the local quadratic cost
\begin{equation}
	\label{eq:local-cost}
	(z_i - s_i)^2 + \sum_{j \in N(i)} w_{i,j} (z_i - z_j)^2 ,
\end{equation}
given the opinions of its neighbors.
Thus, opinion averaging can be viewed as a sequence of best-response updates in a complete-information game, where each agent selects an opinion to balance
conformity with its intrinsic belief.
The resulting Nash equilibrium coincides with the minimizer of the global social cost and is unique due to strong convexity.


\subsection{Exsiting Algorithm} 
  